\documentclass[11pt]{book}

\usepackage{graphicx}
\usepackage{eso-pic}
\usepackage{tcolorbox}
\usepackage{polyglossia}
\setdefaultlanguage[indentfirst=false,forceheadingpunctuation=false]{russian}
\setotherlanguages{english}

\RequirePackage[a4paper, headsep = 0.5 \headsep, left=2.5cm, right=2.1cm, top=2cm, bottom=2.1cm]{geometry}

\setmainfont{Times New Roman}
\newfontfamily\cyrillicfont{Times New Roman}[Script=Cyrillic, Ligatures=TeX]

\usepackage{enumitem}
\setlist[itemize, 1]{label =\raisebox{-0.3\height}{\scalebox{1.8}{\textbullet}}}

\title{Принимая отчужденность}
\author{Тодд МакГоуэн}
\date{2019}
\begin{document}

\AddToShipoutPictureBG*{\includegraphics[width=\paperwidth]{title.png}}

\null
    \thispagestyle{empty}%
    \addtocounter{page}{-1}%
    \newpage

\tableofcontents

\chapter{Введение}

''Отчужденность'' звучит непривлекательно. Будто что-то разрушительное, чего лучше избегать. Никто не говорит друзьям ''Давайте отчуждаться'', наверняка это не так приятно, как основание сообщества, романтические отношения, общение с приятелями, пицца, игра в теннис или выпивка. Напротив --- мы хотим помочь подростку избежать отчужденности, а отчужденному постороннему --- почувствовать себя уютней. Левые, согласно Карлу Марксу, понимают политическую деятельность как борьбу с отчуждением. Почти всем отчужденность видится проблемой, которую надо решить.

Приговор отчужденности поспешен и забывает о ее освободительном качестве. Исторически, у отчужденности плохая репутация, но ведь именно отчужденность дает возможность выхода за пределы конкретной ситуации. Следовало бы думать о ней гораздо лучше. Увидеть ее достоинства и постараться усилить ее опыт.

Никто открыто не отстаивал отчужденность как проект, но некоторые принимают ее основой свободы. Немецкий идеализм, кульминацией которого стала философия Гегеля, говорит об отчужденности как условии самосознания. Психоанализ воспринимает отчужденность субъекта как непреклонный факт нашего существования и стремится примирить с ним. Экзистенциализм Сартра и Симоны де Бовуар тоже видит в отчужденности субъекта основу свободы. Франц Фанон объединяет все три течения --- Гегеля, психоанализ и экзистенциализм --- и формулирует антиколониальную теорию, выступая против попыток избегать фундаментальной нам всем отчужденности. Впрочем, эти мыслители никогда не используют отчужденность как основное понятие или лозунг, хотя оно объединяет эти разрозненные попытки говорить о свободе. В этой книге отчужденность выдвигается на первый план.

Отчужденность парадоксальна. Она не значит отклонения от исходной идентичности, от того, что мы когда-то имели но потеряли. Отчужденность неизменно присуща нам. Она не отдаленность от первичного, потому что первична. Мы отчуждены в своей субъективности до самой смерти. Смерть --- единственное лекарство от этой болезни, так что, наверное, не такая уж это и болезнь. Отчужденность --- не то, о чем можно сожалеть или чего лучше избегать, ведь благодаря ней возможно то, что делает существование достойным продолжения --- свобода, равенство, солидарность, способность выходить за пределы ситуации.

Все творение зависит от отчужденности, она выкорчевывает нас из определенности, освобождает от конкретной ситуации. Каждый рождается в специфичном социуме, в котором на него действуют мощные силы, но не сводится к ним. Первичное отчуждение в языке дистанцирует от других и создает внутреннюю дистанцию: к себе теперь можно относиться как к другому. Мы отчуждены как от самих себя, так и от других, и что бы мы ни делали, нам никогда не преодолеть это отчуждение. Оно основа нашего существования.

Отчужденность --- это отсутствие самоидентичности. Оно отдаляет субъекта от условий, в которых он проявляется. Если я не тождествен себе, если я в противоречии с собой, то внешние силы не могут меня полностью определить. Существо, дистанцированное от самого себя, никогда не сможет быть тем, во что его толкают внешние силы. Этот внутренний раскол значит сопротивление детерминированию биологией или обществом. Именно дистанция от самой себя позволяет субъективности действовать против определяющего. Социальное и натуралистическое утрачивают свою власть на субъектом, потому что его отчужденность делает их незначительными.

Я могу не оставаться до конца своих дней жителем маленького американского городка на Среднем Западе, могу бросить вызов определяющим факторам этой ситуации, и все потому что никогда не отождествлен с положением в котором родился или нахожусь. Это дает мне возможность отнестись к ситуации критически, вплоть до отказа от всех идей и ценностей, которые возложило на меня общество. Я могу быть воспитан с энтузиазмом к капиталистическому строю, и стать марксистом. Если меня заставляют быть женственным, я могу стать мужественным. Решившись полюбить Средний Запад Америки, а, возможно, стану работать в Аргентине, и т.д.. Моя неспособность быть именно тем, кто я есть --- моя не-целостность --- вот что позволяет мне выйти из исходной ситуации. И даже если я остаюсь в ней и принимаю ее, я все равно отношусь к ней как отчужденный. Можно быть чужим и у себя дома. У меня, как и у всех, есть внутренняя чужеродность, мешающая прямому усвоению диктата общества.

Без этой внутренней дистанции человек не был бы субъектом, не смог бы соотноситься с собой. Субъективность порождена и определена этим расколом. Отделенность от самих себя делает нас нами --- неотделимыми от чуждого в нас. Чуждое дает нам нашу уникальность.

Эта внутренняя чуждость и есть то, что Зигмунд Фрейд назовет бессознательным. Бессознательное ставит барьер как между субъектом и ситуацией, так и внутри субъекта. Только отчужденные обладают бессознательным. Вернее, чуждый статус бессознательного значит, что это оно владеет нами. Это иное внутри диктует субъекту его действия, направляет усилия и создает убеждения. Бессознательное ориентирует субъективность, потому что выражает ее самораскол.

\end{document}
